\section{Geometric reduction of the handles}
\label{sec:cancellation}

We continue with the fixed cobordism $(W;A,B)$ and the conventions of
Section~\ref{sec:relative-model}.  The aim is to retain just $b$ handles
in each of indices five and seven.  There are two distinct steps:
elementary operations produce a coefficient exactly equal to $1$, and
Whitney moves turn that coefficient into one geometric intersection.
All moves below occur in the interior of $W$.

\subsection{Introducing the augmentation generators}

A slide of one index-$j$ handle over another replaces a basis vector
$e_j$ by $e_j+e_iq$ in the associated free right $\Lambda$-module.
To realize this, first take $q=\pm g$, with $g\in G$: choose a band
whose path represents $g$, and choose its orientation to give the sign.
The slide arcs can be chosen embedded and disjoint from the other
attaching cores in an eleven-dimensional level; their tubular
neighborhoods can then be shrunk as needed.  Repeated slides
realize an arbitrary finite integral combination of group elements.
The framing is carried through the band construction.  This is the
usual labelled handle slide; see \cite[proof of Theorem~4.2,
pp.~9--10]{Morgan2020}.  The two adjacent differentials change together
as described in Section~\ref{sec:relative-model}.

Write $e_1,\ldots,e_s$ for the existing index-five handles and
$a_i=\partial_1e_i\in I$.  By Lemma~\ref{lem:handle-complex}, the $a_i$
generate $I$ as a right ideal.  Fix a generating list
$y_1,\ldots,y_b$ for $I$, padded with zeros if necessary.

Introduce $b$ cancelling pairs of indices five and six, with new lower
handles $n_1,\ldots,n_b$ satisfying $\partial_1n_j=0$.
Here is a relative placement of the births.  In the level just after
the index-four handle, choose disjoint small balls missing its belt
sphere and all old index-five attaching cores.  These cores and the
belt sphere have dimensions four and seven, respectively, so their
finite union has empty interior.  After choosing the balls, shrink the
old attaching neighborhoods to miss them.  The balls then survive the
old index-five attachments.  Their product traces support ordinary
local five/six births.  When ordering the resulting handles by index,
pull the new lower attaching spheres back along these product traces.
They still miss the index-four belt sphere, so their boundary columns
are zero as asserted.

Choose coefficients in $\Lambda$ such that
\[
 y_j=\sum_{i=1}^s a_iq_{ij},
 \qquad
 a_i=\sum_{j=1}^b y_jt_{ji}.
\]
First slide each new handle to replace $n_j$ by
$n_j+\sum_i e_iq_{ij}$.  Its new column is $y_j$.
Next replace each old handle by $e_i-\sum_j n_jt_{ji}$, using the
new bases.  Every old column is now zero.  The new $b$ columns still
generate $I$.

\subsection{Producing a cancellable pair}

\noindent\emph{A coefficient equal to one.}
Select an old index-five handle $e$, and place it last among the
index-five handles.  Since $\partial_1e=0$, exactness in
Lemma~\ref{lem:handle-complex} gives $e\in\operatorname{im}\partial_2$.
If $h_1,\ldots,h_m$ are the index-six handles and $x_j$ is the
$e$-coefficient in $\partial_2h_j$, it follows that
\[
 x_1\Lambda+\cdots+x_m\Lambda=\Lambda.
\]
While an old handle remains, $r_1\geq b+1$; hence
$m=r_2\geq r_1\geq b+1\geq3$ by \eqref{eq:handle-identities}.
Lemma~\ref{lem:row} therefore applies.  Slides among the index-six
handles produce one whose coefficient into $e$ is exactly $1$.
Place this handle first among the index-six handles, immediately
after $e$.

\smallskip
\noindent\emph{Removing the extra intersections.}
Let $L$ be the intervening regular level, $T\cong S^5$ the upper
handle's attaching core sphere, and $D\cong S^6$ the lower handle's
belt sphere.  Their signed, group-labelled intersection is $1$.
For each group label, pair intersection points with opposite signs,
leaving exactly one positive point with identity label.  This pairing
is possible because all nonidentity coefficients vanish and the
identity coefficient is one.  In particular, the argument uses the
specific coefficient $1$, rather than an arbitrary unit of $\Lambda$.

For each pair, choose arcs joining its points on $T$ and on $D$.
The two spheres are simply connected.  Equality of the group labels
means that the resulting Whitney loop is null-homotopic in $W$.
By \eqref{eq:level-fundamental-group}, it is null-homotopic in $L$.
It therefore bounds a disk in the interior of $L$.  General position
makes this disk embedded, with interior disjoint from $T$, $D$,
the other disks, and the other attaching or belt cores that must be
retained.  The relevant inequalities are
\[
 2+5<11,\qquad 2+6<11,\qquad 2+2<11,
 \qquad 2+7<11.
\]
The last inequality covers the largest other core that occurs in
our index range.  Boundary arcs can likewise be chosen to miss the
other cores and intersection points.  All of these are finite,
compact constructions in $\operatorname{int}L$.

The high-dimensional Whitney theorem now removes each paired
intersection by an isotopy of $T$.  Its hypotheses are complementary
sphere dimensions at least three, opposite intersection signs, and
the null-homotopic Whitney loop just established; simple connectedness
of $L$ is not required.  The framing step is part of this theorem:
the normal splitting along the boundary of the Whitney disk extends
across it.  In the present dimensions the frame-extension space is
$V_4(\R^9)=O(9)/O(5)$, which is simply connected.
See \cite[Theorem~6.6 and Lemma~6.13]{Milnor1965}, or the group-ring
formulation in \cite[Theorem~7.27(i) and Corollary~7.30]{Ranicki2002}.
Transport the existing normal framing of $T$ through its isotopy.
We do not require that this entire handle framing extend over the
Whitney disk.

For support, first arrange the Whitney disks and their compact move
neighborhoods to miss the other cores; then shrink the other tubular
attaching neighborhoods to miss those supports.  A dimension argument
alone would not give disjointness from preassigned full-dimensional
neighborhoods.  The selected belt sphere remains the reference sphere
against which intersections are counted; the isotopy of the attaching
sphere removes intersections with it.  After the finite sequence of
Whitney moves, $T$ meets $D$ transversely in exactly one point.

\smallskip
\noindent\emph{Cancellation relative to the boundary.}
The geometric handle cancellation criterion now applies
\cite[Lemma~8.24]{Ranicki2002}.  We spell out its relative support
here because the quadratic collars must survive.  Isolate the two
adjacent critical points in a slab containing no others, and use a
descending gradient-like field that is a product on the side collar.
The single transverse intersection gives precisely one connecting
trajectory from the index-six point to the index-five point.  Choose
a regular level $c$ between the slab's bottom and the lower critical
value.  Every other nonconstant trajectory in the upper point's unstable manifold
crosses $c$.  Indeed, a trajectory that remained above $c$ could not
escape through the product side; compactness and strict decrease of
the height would force convergence to the only lower critical point,
contrary to uniqueness of the connecting trajectory.

A product side collar has an invariant inner boundary, which the
unstable manifold cannot cross.  Retain a strictly smaller side collar
outside that boundary.  The closure of the unstable manifold above
$c$ then lies in the compact part of the slab between $c$ and the
upper critical value, away from the retained collar and the slab's ends.
This closure can include limiting trajectories through the lower
critical point; it need not be a smooth disk.
The local Morse cancellation theorem gives a deformation supported
in any sufficiently small neighborhood of this closure, fixed elsewhere
\cite[Theorem and Corollary~1]{Laudenbach2013}.
One can apply its closed-manifold formulation by first extending the
height and field to a closed enlargement; they agree on the slab,
and the deformation restricts to the same $W$.
Thus the cancellation changes neither the earlier handles nor the
side and end collars.  The upper attaching maps are transferred
through the induced identification of levels.

For clarity, the chosen index-six column may also have coefficients
into other lower handles.  In bases beginning with the selected pair,
its differential has the block form
\[
 \partial_2=
 \begin{pmatrix}1&a\\ v&B\end{pmatrix},
 \qquad \partial_1=(0\ \ d),\qquad dv=0.
\]
Cancellation leaves $B-va$ as the next differential and leaves $d$
unchanged.  This records algebraically why one coefficient equal to
one suffices and why the other old zero columns of $\partial_1$
survive.  Geometrically this is the preservation of the lower handles
and transfer of the upper attachments just described.

\smallskip
\noindent\emph{Iteration.}
We can therefore repeat the construction until every old index-five
handle has been removed.  Each step cancels one handle of index five
and one of index six.  The identities of Lemma~\ref{lem:handle-complex}
continue to hold for the same cobordism, and the $b$ new lower handles
remain.  At the end,
\[
 r_1=b.
\]

\subsection{The reverse reduction and the quadratic collars}

Turn the cobordism around.  Near its faces the new height is $-Q$,
and an original handle of index $4+i$ has reversed index $8-i$.
The relative covering homology for $(W,B)$ is the same as for $(W,A)$,
so the preceding construction applies without changing the dimension
or adding any variables.  It reduces the reversed index-five count
to $b$, that is, the original index-seven count to $b$.
Reversed five/six births and cancellations are original seven/six
births and cancellations.  They preserve the number $r_1=b$ of original
index-five handles; the upper attaching data in the reversed
presentation are transferred as needed.  Both endpoint counts remain
one.  Lemma~\ref{lem:handle-complex} now gives
\begin{equation}
 (r_0,r_1,r_2,r_3,r_4)
   =\bigl(1,b,\chi(X)-2+2b,b,1\bigr).
 \label{eq:final-handle-counts}
\end{equation}

The handle traces and their intervening product regions supply an
actual smooth Morse height on the same $W$.  Since only finitely many
compactly supported moves were used, a common smaller side collar
still has height $Q$.  We may also prescribe height $-1$ and $1$ on
the two end faces and make its inward derivative strictly positive
at the incoming end and strictly negative at the outgoing end.
It remains to ensure equality with $Q$ on whole end collars, compatible
with the side collar.

At $A$, use the inward coordinate $s=Q+1$.  On a sufficiently small
collar the height has the form
\[
 -1+s\,a(y,s),\qquad a(y,s)>0,
\]
and $a=1$ near the side.  Replace it by
\[
 -1+s\bigl[1+\theta(s)(a(y,s)-1)\bigr],
\]
where $0\leq\theta\leq1$, with $\theta=0$ close to $s=0$ and
$\theta=1$ farther inside.
Put $c_0=\min(1,\inf a)>0$.  Its normal derivative is
\[
 1+\theta(a-1)+s\theta\,\partial_sa+s\theta'(a-1).
\]
The first two terms together are at least $c_0$.  Choose the collar
so small that $|s\partial_sa|<c_0/4$, and choose the transition so
slowly in $\log s$ that
$|s\theta'|\,\|a-1\|_\infty<c_0/4$.
For example, a fixed smooth transition applied to
$\log(s/s_0)/L$, with $L$ large and then $s_0$ sufficiently small,
has these properties and is constant near both endpoints.
The normal derivative is then at least $c_0/2$.
This introduces no critical points and leaves the side collar fixed.
Apply the identical argument at $B$ to one minus the height with
inward coordinate $1-Q$.

The final height equals $Q$ near every face and corner of $W$.
Extend it by $Q$ to $X\times\R^8$ and call the result $F$.
This extension is smooth across all seams.  Its exterior has no
critical points because $dQ=0$ only on the zero section, which lies
in the interior of $W$.  It equals $Q$ for $r\geq3$ and, by
\eqref{eq:final-handle-counts}, has exactly
\[
 1+b+\bigl(\chi(X)-2+2b\bigr)+b+1=\chi(X)+4b
\]
critical points.  This completes the proof of
Proposition~\ref{prop:relative-morse}.
